\begin{pastexam}[label={px:ust-110-ee-ds-01}]{臺聯大電機類 110 學年度 第 1 題}
（問答）

A string matching algorithm aims to find the occurrence of a string P within another string S.

(5\%) (a) Use pseudo-code or C code to describe a string matching algorithm that has time complexity of
O(N$_\mathrm{P}$*N$_\mathrm{S}$), where N$_\mathrm{P}$ and N$_\mathrm{S}$ are the lengths of P and S,
respectively.

(5\%) (b) The Knuch-Morris-Pratt (KMP) algorithm is a string matching algorithm designed to reduce
redundant symbol comparisons in string matching. What is the big-O time complexity, in terms of
N$_\mathrm{P}$ and N$_\mathrm{S}$, of the KMP algorithm?
\end{pastexam}

\begin{pastsoln}
答案：(a) 樸素比對，見下面的程式；(b) $O(N_P + N_S)$。 題目的「Knuch」是 Knuth 的筆誤。

\textbf{(a)} 樸素比對（資結書第~27.1~節）：把 $P$ 放在 $S$ 的每一個可能的起點 $i$，從頭逐字比，
全部相同就是一次出現。 下面用 C 寫，字串從 $0$ 開始編號，印出所有出現的位置。
\begin{minted}[fontsize=\small]{c}
#include <stdio.h>

/* 在 S[0..ns-1] 裡找 P[0..np-1]，回傳出現次數 */
int naive_match(const char *S, int ns, const char *P, int np)
{
    int i, j, count = 0;
    for (i = 0; i + np <= ns; i++) {        /* 起點 i = 0..ns-np */
        j = 0;
        while (j < np && S[i + j] == P[j])  /* 逐字比，不同就停 */
            j++;
        if (j == np) {                      /* np 個字都相同 */
            printf("P occurs at %d\n", i);
            count++;
        }
    }
    return count;
}
\end{minted}
例如 $S = \texttt{abcabcab}$、$P = \texttt{cab}$，起點 $i = 2$ 和 $i = 5$ 都比對成功。

\textbf{正確性。} $P$ 出現在 $S$ 的位置 $i$，意思就是 $S[i + j] = P[j]$ 對 $0 \le j < N_P$ 都成立，
而且 $i + N_P \le N_S$。 外層迴圈試遍所有滿足 $i + N_P \le N_S$ 的 $i$，內層迴圈在第一個不相同的
字停下；$j$ 走到 $N_P$ 表示每個字都相同。 所以每個出現位置都會被印出，印出的也都是出現位置。

\textbf{複雜度。} 起點有 $N_S - N_P + 1$ 個，每個起點最多比 $N_P$ 個字，最壞
$(N_S - N_P + 1) N_P = O(N_P N_S)$ 次比較。 最壞的例子是 $S = \texttt{aa}\cdots\texttt{a}$、
$P = \texttt{a}\cdots\texttt{ab}$：每個起點都要比到 $P$ 的最後一個字才失敗。

\textbf{(b)} KMP 的時間是 $O(N_P + N_S)$：先花 $O(N_P)$ 算 $P$ 的失敗函數（前綴函數），比對時再花
$O(N_S)$（資結書定理~27.4）。 理由是比對時 $S$ 上的指標從不往回走：比對成功時 $S$ 的指標前進一格；
失敗時只依失敗函數把 $P$ 的指標往回跳，$S$ 不動。 $P$ 的指標每次只會加一，往回跳的總次數不會
超過加一的總次數 $N_S$，所以比對階段最多 $2N_S$ 次比較。 失敗函數用同樣的論證是 $O(N_P)$。
它省下的就是 (a) 裡「失敗後 $S$ 的起點只往右移一格、重新從 $P[0]$ 比起」造成的重複比較。
\end{pastsoln}
